% =====================================================================
%  PRÉAMBULE COMMUN — Carnet d'entraînement, Fiche 2 (Nomenclature)
%  (inséré physiquement dans chaque document pour qu'il soit autonome)
% =====================================================================
\documentclass[11pt,a4paper]{article}

% ---------- Encodage & langue ----------
\usepackage[utf8]{inputenc}
\usepackage[T1]{fontenc}
\usepackage{lmodern}
\usepackage[french]{babel}

% ---------- Mathématiques & chimie ----------
\usepackage{amsmath,amssymb,mathtools}
\usepackage{icomma}
\usepackage[version=4]{mhchem}
\usepackage{siunitx}
\sisetup{output-decimal-marker={,},per-mode=symbol}

% ---------- Graphismes & mise en page ----------
\usepackage{xcolor}
\usepackage{colortbl}
\usepackage{tikz}
\usepackage[most]{tcolorbox}
\usepackage{enumitem}
\usepackage{array}
\usepackage{booktabs}
\usepackage{multicol}
\usepackage{fancyhdr}
\usepackage{titlesec}
\usepackage[hidelinks]{hyperref}
\usetikzlibrary{arrows.meta,positioning,calc,shapes.geometric,shapes.misc,
  fit,backgrounds,decorations.pathreplacing}
\usepackage[a4paper,margin=2.2cm,headheight=28pt,headsep=10pt]{geometry}

% =====================================================================
%  PALETTE (charte « Chimie » — mauve/prune du cours)
% =====================================================================
\definecolor{primary}{RGB}{146,95,161}
\definecolor{primarydark}{RGB}{92,55,108}
\definecolor{defcol}{RGB}{0,114,178}
\definecolor{thmcol}{RGB}{0,140,120}
\definecolor{formulacol}{RGB}{198,93,9}
\definecolor{remarkcol}{RGB}{120,94,200}
\definecolor{attncol}{RGB}{200,40,55}
\definecolor{excol}{RGB}{0,135,90}
\definecolor{methcol}{RGB}{90,90,100}
\definecolor{metalcol}{RGB}{198,150,40}
\definecolor{nonmetcol}{RGB}{0,140,150}
\definecolor{oxycol}{RGB}{210,55,45}
\definecolor{hydcol}{RGB}{60,120,210}
\definecolor{catcol}{RGB}{200,70,120}
\definecolor{ancol}{RGB}{40,110,190}
\definecolor{saltcol}{RGB}{120,90,160}

% =====================================================================
%  STYLE COMMUN DES BOÎTES
% =====================================================================
\tcbset{
  boxbase/.style={enhanced,breakable,boxrule=0.9pt,arc=2.5pt,
     left=3mm,right=3mm,top=2.3mm,bottom=2.3mm,
     fonttitle=\bfseries,coltitle=white,
     toptitle=0.6mm,bottomtitle=0.6mm},
}

\newtcolorbox{definition}[1][]{%
  boxbase,colback=defcol!5,colframe=defcol,
  title={\textsc{Définition}\ifx\relax#1\relax\relax\else\textmd{ --- \itshape #1}\fi}}
\newtcolorbox{regle}[1][]{%
  boxbase,colback=thmcol!6,colframe=thmcol,
  title={\textsc{Règle}\ifx\relax#1\relax\relax\else\textmd{ --- \itshape #1}\fi}}
\newtcolorbox{formule}[1][]{%
  boxbase,colback=formulacol!6,colframe=formulacol,
  title={\textsc{Méthode-clé}\ifx\relax#1\relax\relax\else\textmd{ --- \itshape #1}\fi}}
\newtcolorbox{remarque}[1][]{%
  boxbase,colback=remarkcol!6,colframe=remarkcol,
  title={\textsc{Remarque}\ifx\relax#1\relax\relax\else\textmd{ : \itshape #1}\fi}}
\newtcolorbox{attention}[1][]{%
  boxbase,colback=attncol!6,colframe=attncol,
  title={\textsc{Attention}\ifx\relax#1\relax\relax\else\textmd{ : \itshape #1}\fi}}
\newtcolorbox{exemple}[1][]{%
  boxbase,colback=excol!5,colframe=excol,
  title={\textsc{Exemple}\ifx\relax#1\relax\relax\else\textmd{ : \itshape #1}\fi}}
\newtcolorbox{methode}[1][]{%
  boxbase,colback=methcol!7,colframe=methcol,sharp corners,
  title={\textsc{Méthode}\ifx\relax#1\relax\relax\else\textmd{ : \itshape #1}\fi}}
\newtcolorbox{astuce}[1][]{%
  boxbase,colback=primary!6,colframe=primary,
  title={\textsc{Astuce concours}\ifx\relax#1\relax\relax\else\textmd{ : \itshape #1}\fi}}

% ---- Exercice (numéroté) ----
\newtcolorbox[auto counter]{exercice}[1][]{%
  boxbase,colback=primary!4,colframe=primarydark,
  title={\textsc{Exercice}~\thetcbcounter\ifx\relax#1\relax\relax\else\textmd{ --- \itshape #1}\fi}}

% ---- Corrigé (numéroté) ----
\newtcolorbox[auto counter]{corrige}[1][]{%
  boxbase,colback=excol!4,colframe=excol,
  title={\textsc{Corrigé}~\thetcbcounter\ifx\relax#1\relax\relax\else\textmd{ --- \itshape #1}\fi}}

% =====================================================================
%  EN-TÊTES ET PIEDS DE PAGE
% =====================================================================
\pagestyle{fancy}
\fancyhf{}
\fancyhead[L]{\small\textcolor{primarydark}{\textbf{Fiche 2} — Nomenclature}}
\fancyhead[R]{\small\textcolor{primarydark}{\textsc{Carnet d'entraînement}}}
\fancyfoot[C]{\small\thepage}
\renewcommand{\headrulewidth}{0.8pt}
\renewcommand{\headrule}{\hbox to\headwidth{\color{primary}\leaders\hrule height \headrulewidth\hfill}}

\titleformat{\section}{\Large\bfseries\color{primarydark}}{\thesection}{0.6em}{}
\titleformat{\subsection}{\large\bfseries\color{primary}}{\thesubsection}{0.6em}{}
\titleformat{\subsubsection}{\normalsize\bfseries\color{primary!80!black}}{}{0em}{}

% ---- Bandeau de titre réutilisable : \bandeau{Thème}{Sous-titre} ----
\newcommand{\bandeau}[2]{%
  \begin{center}
  \begin{tikzpicture}
  \node[rectangle,rounded corners=7pt,fill=primary,text=white,
        minimum width=16.0cm,minimum height=1.7cm,align=center,inner sep=8pt]
    {{\Large\bfseries #1}\\[3pt]{\normalsize #2}};
  \end{tikzpicture}
  \end{center}
  \vspace{0.3cm}}
\begin{document}
\bandeau{Thème 2 — Du nom à la formule}{Tuto : la neutralité électrique et le croisement des valences}

\noindent Écrire une formule brute à partir d'un nom, c'est appliquer \textbf{une seule} idée : la substance
est \emph{électriquement neutre}. Les charges $+$ doivent compenser exactement les charges $-$. Cette
contrainte fixe les indices (les petits chiffres) de la formule.

\section*{1. La méthode en 4 étapes}
\begin{methode}[passer d'un nom à une formule]
\begin{enumerate}[label=\arabic*.,leftmargin=2.0em,topsep=2pt,itemsep=1pt]
  \item \textbf{Identifier les deux ions} (cation puis anion) à partir du nom.
  \item \textbf{Écrire leur valence} (valeur absolue de la charge).
  \item \textbf{Croiser les valences} : la valence de l'un devient l'indice de l'autre.
  \item \textbf{Vérifier la neutralité} : (nb cations $\times$ charge) $+$ (nb anions $\times$ charge) $= 0$.
\end{enumerate}
\end{methode}

\begin{formule}[le « chassé-croisé »]
Pour un cation \ce{C^{a+}} (valence $a$) et un anion \ce{A^{b-}} (valence $b$) :
\[ \ce{C^{a+}} + \ce{A^{b-}} \;\longrightarrow\; \ce{C_bA_a}. \]
La valence de l'anion ($b$) devient l'indice du cation ; la valence du cation ($a$) devient l'indice de l'anion.
\end{formule}

\begin{center}
\begin{tikzpicture}[scale=1.0]
  \node[font=\large] (C) at (0,0) {\textcolor{catcol}{$\ce{Al^3+}$}};
  \node[catcol,font=\scriptsize,above=0pt of C] {valence 3};
  \node[font=\large] (A) at (4,0) {\textcolor{ancol}{$\ce{SO4^2-}$}};
  \node[ancol,font=\scriptsize,above=0pt of A] {valence 2};
  \draw[->,>=Stealth,catcol,line width=1pt] (0.55,-0.25) to[bend right=18] (3.7,-0.55);
  \draw[->,>=Stealth,ancol,line width=1pt]  (3.55,-0.25) to[bend left=18] (0.4,-0.55);
  \node[catcol,font=\scriptsize] at (2.0,-0.98) {le 3 descend sur l'anion};
  \node[ancol,font=\scriptsize] at (2.0,0.95) {le 2 descend sur le cation};
  \draw[->,>=Stealth,primary,line width=1.3pt] (5.0,0) -- (6.3,0);
  \node[font=\large] (S) at (7.7,0) {\textcolor{saltcol}{$\ce{Al2(SO4)3}$}};
  \node[saltcol,font=\scriptsize,below=1pt of S] {sulfate d'aluminium};
\end{tikzpicture}
\end{center}

\section*{2. Deux gestes à ne pas oublier}
\begin{regle}[simplifier et mettre des parenthèses]
\begin{itemize}[leftmargin=1.5em,topsep=2pt,itemsep=1pt]
  \item \textbf{Simplifier} les indices par leur diviseur commun. Ex.\ \ce{Ca^2+}$+$\ce{O^2-} donne
        \ce{Ca2O2}, que l'on simplifie en \ce{CaO}.
  \item \textbf{Parenthèses} : tout ion \emph{polyatomique} affecté d'un indice $\geq 2$ se met entre
        parenthèses. Ex.\ \ce{Ca(OH)2}, \ce{Al(NO3)3}, \ce{(NH4)2SO4}. On n'en met jamais pour un indice 1
        (\ce{NaNO3}) ni pour un ion monoatomique (\ce{CaCl2}).
\end{itemize}
\end{regle}

\begin{exemple}[nitrate d'étain(IV)]
\textbf{1)} cation \ce{Sn^4+} (le « (IV) » donne le degré), anion nitrate \ce{NO3-}. \quad
\textbf{2)} valences 4 et 1. \quad
\textbf{3)} croisement : $\ce{Sn^4+} + \ce{NO3-} \to \ce{Sn(NO3)4}$ (parenthèses car \ce{NO3} $\times 4$). \quad
\textbf{4)} neutralité : $(+4) + 4(-1) = 0$. \checkmark \quad
$\boxed{\ce{Sn(NO3)4}}$
\end{exemple}

\begin{attention}
Ne confonds jamais les couples piégeux : nitr\textbf{ite} \ce{NO2-} / nitr\textbf{ate} \ce{NO3-} ;
sulf\textbf{ite} \ce{SO3^2-} / sulf\textbf{ate} \ce{SO4^2-} ; thiosulfate \ce{S2O3^2-} / sulfate.
Une lettre change l'anion, donc la formule.
\end{attention}
\end{document}
